\section{From LLM Reasoning to Agentic Reasoning}
\label{sec:from}

Traditional reasoning with large language models (LLMs) is typically formulated as a one-shot or few-shot prediction task over static inputs. These models rely on scaling \textbf{test-time computation}, improving accuracy by increasing model size or inference budget, but without the ability to interact, remember, or adapt to changing goals. Methods such as prompt engineering, in-context learning, and chain-of-thought prompting have made reasoning more explicit, yet conventional LLMs remain passive sequence predictors that operate within fixed prompts.

\emph{Agentic reasoning}, in contrast, emphasizes \textbf{scaling test-time interaction}. Instead of depending solely on internal parameters, agentic systems reason through action: invoking tools, exploring alternatives, updating memory, and integrating feedback. This transforms inference into an iterative process that includes decision steps, reflection, and learning from experience. Reasoning becomes a dynamic loop that connects the model, memory, and environment.


\begin{table}[h]
\centering
\caption{Contrasting capabilities of \textbf{LLM reasoning} and \textbf{agentic reasoning}.}
\label{tab:llm-vs-agentic}
\renewcommand{\arraystretch}{1.25}
\begin{tabular}{
  >{\centering\arraybackslash\bfseries}m{3.2cm}
  @{\hspace{3pt}}
  >{\centering\arraybackslash}m{4.5cm}
  @{\hspace{2pt}}
  >{\centering\arraybackslash}m{0.5cm}
  @{\hspace{2pt}}
  >{\centering\arraybackslash}m{4.5cm}
}
\toprule
\textbf{Dimension}
& \llmchip{\textbf{LLM Reasoning}}
& $\leftrightarrow$
& \agentchip{\textbf{Agentic Reasoning}} \\
\midrule
\multirow{2}{3.2cm}{\centering Paradigm}
  & \llmchip{passive}
  & $\leftrightarrow$
  & \agentchip{interactive} \\
  & \llmchip{static input}
  & $\leftrightarrow$
  & \agentchip{dynamic context} \\
\midrule
\multirow{2}{3.2cm}{\centering Computation}
  & \llmchip{single pass}
  & $\leftrightarrow$
  & \agentchip{multi step} \\
  & \llmchip{internal compute}
  & $\leftrightarrow$
  & \agentchip{with feedback} \\
\midrule
\multirow{2}{3.2cm}{\centering Statefulness}
  & \llmchip{context window}
  & $\leftrightarrow$
  & \agentchip{external memory} \\
  & \llmchip{no persistence}
  & $\leftrightarrow$
  & \agentchip{state tracking} \\
\midrule
\multirow{2}{3.2cm}{\centering Learning}
  & \llmchip{offline pretraining}
  & $\leftrightarrow$
  & \agentchip{continual improvement} \\
  & \llmchip{fixed knowledge}
  & $\leftrightarrow$
  & \agentchip{self evolving} \\
\midrule
\multirow{2}{3.2cm}{\centering Goal Orientation}
  & \llmchip{prompt based}
  & $\leftrightarrow$
  & \agentchip{explicit goal} \\
  & \llmchip{reactive}
  & $\leftrightarrow$
  & \agentchip{planning} \\
\bottomrule
\end{tabular}
\end{table}







This transition marks a conceptual shift: reasoning no longer scales through static capacity, but through structured interaction that enables planning, adaptation, and collaboration across time and tasks.

\subsection{Positioning Our Survey}

While several recent surveys have examined LLM reasoning or agent architectures \citep{huang2022towards, chen2025towards, xu2025towards, ke2025survey, zhang2025survey, zhang2025landscape, lin2025comprehensive, fang2025comprehensive, gao2025survey}, our work focuses specifically on \textbf{agentic reasoning} as a unified paradigm for understanding reasoning as interaction. We position this survey at the intersection of model-centric reasoning and system-level intelligence, aiming to bridge prior discussions on reasoning mechanisms and agent architectures.

\textbf{Relation to LLM Reasoning Surveys.}
Existing surveys on LLM reasoning mainly investigate how to elicit or enhance reasoning within a model’s internal computation process.
For example, \citet{huang2022towards, chen2025towards, xu2025towards, ke2025survey} summarize prompting and scaling techniques such as chain-of-thought, reinforcement post-training, and long-context reasoning, emphasizing how LLMs can learn to reason better through inference-time supervision or post-training alignment.
These works improve the internal expressiveness of reasoning traces but typically remain within static inference settings, where reasoning unfolds in a single forward pass without external interaction.
In contrast, our survey examines how reasoning extends \emph{beyond} text generation, encompassing dynamic planning, adaptive memory, and feedback-driven behavior during deployment.

\textbf{Relation to AI Agent Surveys.}
Several contemporary surveys have begun to explore LLM-based agents from architectural or system perspectives \citep{zhang2025landscape, lin2025comprehensive, fang2025comprehensive, gao2025survey}.
These works analyze how agents employ reinforcement learning, planning, and tool-use modules to operate in complex environments.
For instance, \citet{zhang2025landscape, lin2025comprehensive} focus on reinforcement learning for agentic search and decision-making, while \citet{fang2025comprehensive, gao2025survey} emphasize self-evolving and lifelong agentic systems that continuously learn from interaction.
Our focus complements these perspectives by centering on the \textbf{reasoning process} that these architectures enable, specifically how interaction, feedback, and collaboration transform static inference into adaptive reasoning.
Rather than viewing reasoning as an implicit by-product of architectural design, we treat it as the unifying mechanism that links single-agent reinforcement, multi-agent coordination, and self-evolving intelligence.

In summary, our survey provides a reasoning-centric lens on intelligent agency.
We examine how foundational reasoning mechanisms, post-training adaptation, and long-term self-evolution jointly constitute the basis of \emph{agentic reasoning}, illustrating the transition from static prediction to interactive, adaptive, and continually improving intelligence.


\subsection{Preliminaries}
\label{sec:prelims}

This subsection formalizes the transition from static language modeling to agentic reasoning. To align with the \textbf{three-layered dimensions} (Foundational, Self-Evolving, Collaboration) outlined in the introduction, we unify these capabilities under a single {control-theoretic framework}.

\paragraph{Formalizing Agentic Reasoning: A Latent-Space View.}
Standard approaches often conflate the agent's context with the environment state. We model the environment as a \textbf{Partially Observable Markov Decision Process (POMDP)} and introduce an internal \emph{reasoning variable} to expose the ``think--act'' structure of agentic policies. Concretely, we consider the tuple
$\langle \mathcal{X}, \mathcal{O}, \mathcal{A}, \mathcal{Z}, \mathcal{M}, \mathcal{T}, \Omega, \mathcal{R}, \gamma \rangle$, where
$\mathcal{X}$ is the latent \emph{environment state} space (unobservable to the agent),
$\mathcal{O}$ is the observation space (e.g., user queries, API returns),
$\mathcal{A}$ is the external action space (e.g., tool invocation, final answer),
$\mathcal{Z}$ is a \emph{reasoning trace} space (e.g., latent plans, optionally verbalized as chain-of-thought),
and $\mathcal{M}$ is the agent's \emph{internal memory/context} space (e.g., a sufficient statistic of interaction history).
$\mathcal{T}$ and $\Omega$ denote the transition and observation kernels, $\mathcal{R}$ the reward, and $\gamma\in(0,1)$ the discount factor.

At timestep $t$, the agent conditions on a history
$h_t = (o_{\le t}, z_{<t}, a_{<t})$ (i.e., $o_t$ is observed before generating $z_t$ and then $a_t$).
Equivalently, the history can be summarized by an internal memory state $m_t\in\mathcal{M}$.
Crucially, we distinguish external actions from internal reasoning. We factorize the policy as
\begin{equation}
\label{eq:policy_decomp}
\pi_\theta(z_t, a_t \mid h_t)
= \underbrace{\pi_{\text{reason}}(z_t \mid h_t)}_{\text{Internal Thought}}
\cdot \underbrace{\pi_{\text{exec}}(a_t \mid h_t, z_t)}_{\text{External Action}}.
\end{equation}
This decomposition highlights the core shift in agentic systems: performing computation in $\mathcal{Z}$ (thinking) before committing to $\mathcal{A}$ (acting). The objective remains maximizing the expected return
$J(\theta)=\mathbb{E}_{\tau}\!\left[\sum_{t\ge 0}\gamma^t r_t\right]$.

\paragraph{In-Context Reasoning: Inference-Time Search.}
In this regime, model parameters $\theta$ are frozen. The agent optimizes the reasoning trajectory by searching over $\mathcal{Z}$ to maximize a heuristic value function $\hat{v}(h_t, z)$.
We model inference as selecting a trajectory $\tau = (h_0, z_0, a_0, h_1, z_1, a_1, \ldots)$.
Methods like ReAct \citep{yao2023react} perform greedy decoding over alternating thoughts $z$ and actions $a$.
Tree-of-Thoughts (ToT \citep{yao2023tree}) and related MCTS-style approaches treat partial thoughts as nodes $u \in \mathcal{U}$ (e.g., a representation derived from $(h_t,z_t)$) and search for an optimal path:
\begin{equation}
\tau^\star \in \arg\max_{\tau} \sum_{t} \hat{v}_\phi(u_t),
\end{equation}
where $\hat{v}_\phi$ is a heuristic evaluator or verifier. This corresponds to planning in $\mathcal{Z}$ without updating the policy parameters.

\paragraph{Post-Training: Policy Optimization.}
This paradigm optimizes $\theta$ to align the policy with long-horizon rewards $r_t$ (e.g., correctness, safety), including reasoning models (e.g., DeepSeek-R1 \citep{guo2025deepseekr1}) and learning-to-search systems (e.g., Search-R1 \citep{jin2025search}, DeepRetrieval \citep{jiang2025deepretrieval}) that train multi-turn reasoning or tool use with RL.
While PPO \citep{schulman2017proximal} is standard, \textbf{Group Relative Policy Optimization (GRPO)} \citep{shao2024deepseekmath}-based methods are widely used for reasoning tasks. GRPO eliminates the value network by constructing advantages from group-relative rewards.
For a group of $G$ sampled outputs $\{y_i\}_{i=1}^G$ from the same prompt $q$, a common GRPO objective is:
\begin{equation}
\mathcal{L}^{\text{GRPO}}(\theta)
= \mathbb{E}_{q \sim P(Q)} \left[
\frac{1}{G} \sum_{i=1}^G
\left(
\min\!\left( \rho_i \hat{A}_i,\; \mathrm{clip}(\rho_i, 1-\epsilon, 1+\epsilon)\hat{A}_i \right)
- \beta\, \mathbb{D}_{KL}(\pi_\theta \,\|\, \pi_{\text{ref}})
\right)
\right],
\end{equation}
where $\rho_i = \frac{\pi_\theta(y_i \mid q)}{\pi_{\theta_{\text{old}}}(y_i \mid q)}$ and the group-normalized advantage is
\begin{equation}
\hat{A}_i
= \frac{r_i - \mu}{\sigma + \delta},\quad
\mu=\frac{1}{G}\sum_{j=1}^G r_j,\quad
\sigma=\sqrt{\frac{1}{G}\sum_{j=1}^G (r_j-\mu)^2},
\end{equation}
with $\delta>0$ a small constant for numerical stability.
Advanced methods such as ARPO \citep{lu2025arpo} and DAPO \citep{yu2025dapo} extend this framework to handle sparse rewards and improve stability in complex tool-use environments (e.g., via replay/rollout strategies and decoupled clipping).

\paragraph{Collective Intelligence: Multi-Agent Reasoning.}
We extend the single-agent formulation to a \emph{decentralized} partially observable multi-agent setting, commonly formalized as a \emph{Dec-POMDP}. The core distinction lies in expanding each agent's observation to include a \textbf{communication channel} $\mathcal{C}$. For a system of $N$ agents, the joint policy $\boldsymbol{\pi}$ is composed of individual policies $\pi^i$, where agent $i$'s observation $o^i_t$ explicitly includes communicative messages $c^{-i}_{t-1}$ generated by peers.
Crucially, in agentic MARL, communication is not merely signal transmission but an extension of the reasoning process: one agent's external action can act as a prompt that triggers another agent's internal reasoning chain.
Existing frameworks like AutoGen \citep{wu2024autogen} and CAMEL \citep{li2023camel} represent static role-playing with fixed policies. Recent agentic RL advances (e.g., GPTSwarm~\cite{zhuge2024gptswarm}, MaAS, agents trained via PPO/GRPO~\cite{hong2025multi}) aim to \emph{optimize} this joint reasoning distribution. The challenge shifts from single-agent planning to \textbf{mechanism design}: optimizing the communication topology and incentive structures to align decentralized reasoning processes $\pi^i_{\text{reason}}$ toward a coherent global objective, often utilizing Centralized-Training/Decentralized-Execution (CTDE) paradigms to stabilize the emergence of cooperative behaviors.

\paragraph{Self-Evolving Agents: The Meta-Learning Loop.}
While foundational agents optimize reasoning $z$ within an episode, self-evolving agents optimize the agent system itself across episodes $k=1, \dots, K$. Let $\mathcal{S}_k$ denote the evolvable system state (e.g., explicit memories, tool libraries, or code). A generic meta-update rule is
\begin{equation}
\mathcal{S}_{k+1} \leftarrow U(\mathcal{S}_k, \tau_k, \mathcal{F}_k),
\end{equation}
where $\mathcal{F}_k$ represents environmental feedback (rewards, execution errors) and $\mathcal{S}_k$ represents the evolvable state. We categorize self-evolution by the nature of $\mathcal{S}$:
\begin{itemize}[leftmargin=*]
    \item \textbf{Verbal Evolution:} $\mathcal{S}$ consists of textual reflections or guidelines. Methods like Reflexion \citep{shinn2023reflexion} update $\mathcal{S}$ by synthesizing error logs into linguistic cues that condition future reasoning policies.
    \item \textbf{Procedural Evolution:} $\mathcal{S}$ consists of a library of executable tools or skills. Agents like Voyager \citep{wang2023voyager} evolve by synthesizing new code-based skills, expanding the action space $\mathcal{A}$ permanently.
    \item \textbf{Structural Evolution:} $\mathcal{S}$ consists of the agent's source code or architecture itself. Advanced methods like AlphaEvolve \citep{novikov2025alphaevolve} treat the agent's code as a hypothesis space, using an LLM as a mutation operator to search for superior reasoning algorithms.
\end{itemize}

This framework unifies these diverse approaches as gradient-free or gradient-based optimization steps over the agent's explicit memories and artifacts (and optionally parameters), closing the loop between experience and competence.
